\section{Digital, barrier and Asian options}

Three common exotic payoffs under the lognormal model: a digital pays a fixed amount or the asset on a condition at expiry, a barrier option is switched on or off by the path's level, and an Asian option pays on an average. Their closed forms assume constant volatility and continuous monitoring or averaging; the examples show what changes when monitoring and averaging are discrete. Pricing on grids and by simulation is in \hyperref[meth:finite-difference-option-pricing]{finite differences} and \hyperref[meth:monte-carlo-simulation]{Monte Carlo simulation}.

\begin{qfmethod}[Digital (binary) options: cash-or-nothing and asset-or-nothing values, the vanilla call as their difference, call-spread replication and pin risk]\label{meth:digital-option-valuation}
\qffield{Statement} A cash-or-nothing call pays a fixed amount $Q$ if the asset ends above the strike and nothing otherwise; since the risk-neutral probability of finishing above $K$ is $N(d_2)$, its value is $Qe^{-rT}N(d_2)$. An asset-or-nothing call pays the asset price on the same event and is worth $S_0e^{-qT}N(d_1)$. A regular European call is a long asset-or-nothing call and a short cash-or-nothing call paying $K$; the put decomposes the same way with the puts \citep[chap.~26, sec.~26.10, pp.~622--623]{s7-hulloptions}. Wilmott derives the binary call value $e^{-r(T-t)}N(d_2)$ from the risk-neutral density, and binary call plus binary put equals $e^{-r(T-t)}$ \citep[chap.~7, pp.~119--121; chap.~2, pp.~43--44]{wide-wilmott-pwoqf}. Discontinuous payoffs invite manipulation when the underlying is thin: the holder of a cash digital one day from expiry with the price just below the strike may buy the asset to lift it over (Hull, p.~623).

\qffield{Notation} $d_{1,2}=[\ln(S_0/K)+(r-q\pm\tfrac12\sigma^2)T]/(\sigma\sqrt T)$; unit cash digital $D_c=e^{-rT}N(d_2)$; asset digital $D_a=S_0e^{-qT}N(d_1)$; so $c=D_a-KD_c$. Interpretation: the centred call spread $[c(K-\varepsilon)-c(K+\varepsilon)]/(2\varepsilon)$ tends to $D_c$ with error of order $\varepsilon^2$, because it is a central difference of $-\partial c/\partial K=D_c$. The delta of the unit cash digital, $e^{-rT}n(d_2)/(S\sigma\sqrt T)$, grows like $1/\sqrt T$ at the strike as expiry approaches: the hedge of the step payoff becomes unbounded there, which is pin risk. Because $D_c=-\partial c/\partial K$ follows from the payoffs alone, a smile changes the digital's value: by the chain rule, $-\frac{d}{dK}c(K,\sigma(K))=D_c-\mathcal V\,\partial\sigma/\partial K$, with $\mathcal V$ the call's vega.

\qffield{Conditions} Lognormal prices with constant $r$, $q$ and $\sigma$. Under a smile, a flat volatility misprices the digital by the vega-times-skew term above (\hyperref[def:implied-volatility-surface]{volatility surface}). A call spread of finite width approximates the digital with a bounded delta; the closer the strikes, the closer the value and the larger the delta.

\qffield{Practice} Example, synthetic: $S_0=K=100$, $r=5$ percent, $q=2$ percent, $\sigma=25$ percent, $T=1$. The call is 11.1238, the unit cash digital 0.47372 and the asset digital 58.4955. Then $58.4955-100\times0.47372=11.1238$, equal to $10^{-12}$. Call spreads of half-width 2, 1 and 0.5 miss the digital by $1.03\times10^{-4}$, $2.58\times10^{-5}$ and $6.45\times10^{-6}$, each a quarter of the last. A simulation of 100{,}000 terminal prices gives 0.4719 (standard error 0.0015). At the strike the digital's delta is 0.015 a year before expiry, 0.055 a month before and 0.253 a day before. A digital paying 1~million then needs about 253{,}000 shares, 25~million of stock. The delta stays positive on both sides of the strike (0.2530 at 99.9, 0.2533 at 100 and 0.2525 at 100.1), so the hedge does not change sign; what makes it hard is its size, which falls off steeply as the price moves away from the strike and goes to zero at expiry, while the gamma changes sign near the strike. A price that oscillates around the strike forces large purchases and sales of stock.

\qffield{Pitfalls} Pricing a digital with the at-the-money volatility while ignoring the skew. Delta hedging a short digital into expiry near the strike. Confusing $N(d_2)$, a risk-neutral probability, with the real-world probability of finishing in the money (\hyperref[res:risk-neutral-pricing-feynman-kac]{risk-neutral pricing}).

\qffield{See also} \hyperref[mod:black-scholes-merton]{BSM}; \hyperref[res:put-call-forward-parity]{put--call parity}; \hyperref[meth:option-spreads-volatility-strategies]{option spreads}; \hyperref[meth:barrier-option-valuation]{barrier options}.

\end{qfmethod}

\begin{qfmethod}[Barrier options: knock-in and knock-out, in--out parity, the down-and-out call in closed form, discrete monitoring and the hedging problems near the barrier]\label{meth:barrier-option-valuation}
\qffield{Statement} A knock-out option ceases to exist when the asset reaches a barrier $H$; a knock-in option comes into existence only if it does. Because a regular call is a down-and-in plus a down-and-out call, $c_{do}=c-c_{di}$. For $H\le K$,
\[ c_{di}=S_0e^{-qT}(H/S_0)^{2\lambda}N(y)-Ke^{-rT}(H/S_0)^{2\lambda-2}N(y-\sigma\sqrt T),\quad \lambda=\frac{r-q+\sigma^2/2}{\sigma^2},\quad y=\frac{\ln[H^2/(S_0K)]}{\sigma\sqrt T}+\lambda\sigma\sqrt T, \]
with the other cases given similarly. These formulas assume continuous observation of the price. For observation $m$ times, Broadie, Glasserman and Kou replace $H$ by $He^{0.5826\sigma\sqrt{T/m}}$ for an up barrier and $He^{-0.5826\sigma\sqrt{T/m}}$ for a down barrier. Vega can be negative: for an up-and-out call near the barrier, higher volatility raises the chance of knocking out \citep[chap.~26, sec.~26.9, pp.~620--622]{s7-hulloptions}. Barrier options have a discontinuous delta at the barrier, so ``delta hedging through the barrier is virtually impossible, and certainly very costly''. When gamma changes sign, a single vega number can mislead, and practitioners price with volatilities from vanillas and one-touch options or with a calibrated surface \citep[chap.~23, secs.~23.7--23.8, pp.~401--405]{wide-wilmott-pwoqf}.

\qffield{Notation} $c_{do}$, $c_{di}$ down-and-out and down-and-in calls; $m$ monitoring dates. Interpretation: between two monitoring dates a Brownian path that starts at $x_0=\ln(S_i/H)>0$ and ends at $x_1>0$ stays above the barrier with probability $1-\exp[-2x_0x_1/(\sigma^2\Delta t)]$. Weighting each simulated payoff by the product of these probabilities gives an unbiased estimate of the continuously monitored price from a finite grid.

\qffield{Conditions} Lognormal prices, constant parameters, no rebate. The shift is an asymptotic correction for many monitoring dates, not an exact formula. In a tree or grid, a barrier that falls between nodes causes errors of order $\sqrt{\Delta t}$; interpolating between inner and outer barriers, or placing nodes on the barrier, reduces them \citep[chap.~27, sec.~27.6, pp.~656--658]{s7-hulloptions}. Spikes in thin markets can knock an option in or out (Hull, pp.~622--623).

\qffield{Practice} Examples, synthetic, with the inputs of the digital example and a barrier of 90 (below $S_0$ and $K$):
\begin{itemize}
\item \emph{Closed form.} $c_{do}=8.1388$ and $c_{di}=2.9850$, summing to the call 11.1238. On 100{,}000 daily paths, knock-in plus knock-out equals the vanilla payoff on every path to $10^{-12}$.
\item \emph{Continuous monitoring.} Simulation with Brownian-bridge weights on 252 steps gives 8.1953 (standard error 0.053), within 1.1 standard errors.
\item \emph{Daily monitoring.} Simulation gives 8.4890 (standard error 0.054) against 8.5190 from the closed form with the barrier shifted to 89.178, inside the stated three standard errors plus 0.5 percent. Monitoring daily rather than continuously adds about 0.35.
\item \emph{Vega.} For this down-and-out call, vega falls from 14.8 at $S=130$ to 0.8 at $S=91$. An up-and-out call with $K=100$, $H=120$ and $S=115$ is worth 0.20 and has a vega of $-2.2$.
\end{itemize}

\qffield{Pitfalls} Pricing a daily-monitored contract with the continuous formula: it undervalues a knock-out and overvalues a knock-in. Hedging a short knock-out near the barrier with delta alone. Quoting one implied volatility for a barrier whose gamma changes sign.

\qffield{See also} \hyperref[meth:digital-option-valuation]{digital options}; \hyperref[meth:finite-difference-option-pricing]{finite differences}; \hyperref[def:implied-volatility-surface]{volatility surface}.

\qffield{Further reading} \nolinkurl{vo-broadie-glasserman-kou-1997} (the continuity correction and its derivation).

\end{qfmethod}

\begin{qfmethod}[Asian options: average-price payoffs, the geometric average in closed form, the arithmetic average by simulation with a geometric control, and discrete against continuous averaging]\label{meth:asian-option-valuation}
\qffield{Statement} An average-price call pays $\max(0,S_{\text{ave}}-K)$, where $S_{\text{ave}}$ is the arithmetic average of the underlying over the option's life. Such options ``tend to be less expensive than regular options'' and suit a corporate treasurer exposed to an average exchange rate. When the asset follows geometric Brownian motion the geometric average is exactly lognormal and the arithmetic average approximately so. A common approximation fits a lognormal distribution to the first two moments of the arithmetic average and uses Black's model. In Hull's example ($S_0=K=50$, $r=10$ percent, $\sigma=40$ percent, $T=1$) continuous averaging gives 5.62, and 12, 52 and 250 observations give 6.00, 5.70 and 5.63 \citep[chap.~26, sec.~26.13, pp.~626--627]{s7-hulloptions}. For geometric averaging the formula is the vanilla one with volatility $\sigma/\sqrt3$ and an adjusted dividend yield (Kemna and Vorst). The square-root-of-three rule arises because volatility early in the averaging period affects more of the average than volatility late in it \citep[chap.~25, secs.~25.9--25.11, pp.~440--442]{wide-wilmott-pwoqf}. Monte Carlo handles arithmetic averages directly (Hull, chap.~27, sec.~27.5, p.~653).

\qffield{Notation} Averaging dates $t_i=iT/n$, $i=1,\dots,n$, excluding $S_0$; $G=(\prod S_{t_i})^{1/n}$. Interpretation, derived from the lognormal model: $\ln G$ is normal with mean $\ln S_0+(r-q-\tfrac12\sigma^2)T\frac{n+1}{2n}$ and variance $\sigma^2T\frac{(n+1)(2n+1)}{6n^2}$, and the geometric call is $e^{-rT}[e^{m+v/2}N(d_1)-KN(d_2)]$ with $d_2=(m-\ln K)/\sqrt v$ and $d_1=d_2+\sqrt v$. As $n\to\infty$ the variance tends to $\sigma^2T/3$ and the mean to $\ln S_0+\tfrac12(r-q-\tfrac12\sigma^2)T$. This is Black--Scholes with volatility $\sigma/\sqrt3$ and yield $\tfrac12(r+q+\sigma^2/6)$, the form of Wilmott's full formula (p.~442).

\qffield{Conditions} Lognormal prices with constant parameters. The discrete formula must match the contract's averaging dates, including whether $S_0$ counts; the continuous formula is its limit, approached at first order in $1/n$. Wilmott's summary on p.~440 gives the adjusted yield as $D+\sigma^2/6$, which does not agree with his own formula on p.~442; the formula on p.~442 is the one that matches the limit. The geometric option is a good control for the arithmetic one because their payoffs are highly correlated, but the arithmetic price still has no closed form. The arithmetic average is never below the geometric one (AM--GM), so the arithmetic call is worth at least the geometric call; the comparison with the vanilla call has no such general ordering. Averaging lowers the dispersion but also moves the forward level, since $E[S_{\text{ave}}]=S_0\frac1n\sum_ie^{(r-q)t_i}$ under the risk-neutral measure: with $q>r$ the average's forward lies above the terminal forward, and an in-the-money Asian call can be worth more than the vanilla.

\qffield{Practice} Examples, synthetic: the inputs of the digital example, 12 monthly averaging dates, $S_0$ excluded, 100{,}000 paths.
\begin{itemize}
\item \emph{Geometric.} The discrete closed form gives 6.3903 and simulation 6.3888 (standard error 0.031).
\item \emph{Arithmetic.} Uncontrolled simulation gives 6.6797 (standard error 0.032). With the geometric payoff as control (estimated coefficient 1.039), the estimate is 6.6813 (standard error 0.0011): variance falls by a factor of about 900.
\item \emph{Averaging and volatility.} With these inputs both averages lie below the vanilla call, 11.1238: averaging lowers the effective volatility, and with $r>q$ it also lowers the forward of the average. Continuous geometric averaging gives 5.9802, 0.41 below 12 monthly dates; the gap halves when the dates double from 500 to 1{,}000. Using the yield $q+\sigma^2/6$ alone would give 6.5181.
\item \emph{Negative carry.} With $S_0=100$, $K=60$, $r=2$ percent, $q=8$ percent, $\sigma=20$ percent and the same 12 dates, the vanilla call is 33.5611 and the geometric Asian call 35.7601. The arithmetic call is at least $e^{-rT}(E[S_{\text{ave}}]-K)=36.0876$, because $\max(0,x)\ge x$. Both Asian calls are worth more than the vanilla.
\end{itemize}

\qffield{Pitfalls} Valuing a monthly-averaged contract with the continuous formula. Forgetting whether the first fixing is today's price. Assuming that an Asian option is always cheaper than the matching vanilla, when carry moves the average's forward. Reporting a controlled estimate without checking it against the uncontrolled one on the same paths.

\qffield{See also} \hyperref[meth:monte-carlo-simulation]{Monte Carlo simulation} (control variates); \hyperref[meth:black-model-rates-currencies]{Black's model}; \hyperref[meth:digital-option-valuation]{digital options}.

\end{qfmethod}
